% !TeX root = ../main.tex
\begin{chineseabstract}
% 如有基金资助可打开下一行
% \footnotetext{*本研究得到某某基金（编号：）的资助}
% \astfootnote{本研究得到某某基金（编号：）的资助}

随着深度神经网络与大语言模型在云端与边缘端的广泛部署，人工智能对硬件算力的需求飞速增长。神经网络芯片作为领域专用架构被提出，并从单核向多核架构演进。然而，多核芯片在流片之前的端到端性能评估面临两难：硬件模拟平台精度高但成本昂贵、迭代缓慢、参数难以灵活探索。解析模型迭代代价低但难以刻画片上网络拥塞、片外存储争用与多核同步等系统级动态行为。周期级模拟器位于二者之间，是芯片在流片之前进行性能评估的关键基础设施。然而，面向多核神经网络芯片、可直接以调度器中间表示为输入的开源模拟器仍属空白。

针对上述空白，本文设计并实现了一套通过调度器中间表示驱动的面向多核神经网络芯片的周期级模拟器。通过统一任务抽象，模拟器将多核执行刻画为核心级工作负载、显式依赖与数据流模型，在统一时间轴与多时钟同步下驱动计算、互连与存储协同推进。模拟器能够记录多核加载、计算、写回以及网络背压等阶段周期，输出芯片状态统计数据，并以热力图、甘特图的形式将模拟结果可视化，形成从解析到执行再到统计的闭环，使瓶颈归因不再依赖人工经验判断。

实验详细验证了模拟器的精度和设计空间探索价值。在精度验证上，本文以商用的 Synopsys ZeBu 硬件模拟平台的执行结果作为基准数据，在不同工作负载和不同批量大小配置下，模拟器输出的总周期与基准数据偏差不超过 $\pm$5.2\%，达到接近硬件模拟的精度。本文进一步围绕多核神经网络芯片开展系统性瓶颈分析与设计空间探索。在四核测试芯片运行 ResNet50 模型的场景下，本文模拟了从片外内存带宽升级再到片上内存端口加宽的递进式优化路径，将端到端总周期降低 73.7\%（3.80 倍加速），并揭示了“片外带宽瓶颈到片上端口瓶颈”的逐级迁移规律。本文进一步沿单核算力、核数扩展与批量大小三个维度展开 12 组设计空间探索，定量发现以下三点：（1）存储带宽低于计算强度要求时提升算力对吞吐量无贡献；（2）芯片核数超过负载并行度上限后性能将严重退化甚至死锁；（3）受限于存储带宽争用，工作负载的批量大小扩展效率仅约 24\%。综上，本文为多核神经网络芯片在流片之前的系统级瓶颈定位与软硬件协同优化提供了可复现的周期级模拟基础与实用评估能力。

\chineseKeywords{神经网络芯片；多核芯片架构；周期级模拟器；设计空间探索}

% 关键词由3～5个组成。关键词应从《汉语主题词表》中摘选，当《汉语主题词表》的词不足以反映主题时，可由申请人设计关键词，但须加注。每一关键词之间用分号分开，最后一个关键词后不打标点符号。由申请人设计的关键词，须在该关键词的右上角标注*，并在该页的页脚处注明“*表示非汉语主题词”。\newline

\chineseType{应用研究}

% 论文类型包括：a.理论研究(Theoretical Research)；b.应用基础(Application Fundamentals)；c.应用研究(Application Research)；d.研究报告(Research Report)；e.设计报告(Design Report)；f.案例分析(Case Study)；g.调研报告(Investigation Report)；h.产品研发(Product Development)；i.工程设计(Engineering Design)；j.工程/项目管理(Engineering/Project Management)；k.其它（Others）。

\end{chineseabstract}

